\subsection*{From the book to a ready spell: a complete example}
\begin{enumerate}
\item Enter CAMP and choose Memorize Spells for a mage, or Pray for Spells for a cleric or paladin.
\item Choose the character and spell level. Add copies with a left click or the + button; remove them with a right click or the $-$ button. If 0 remain, free a slot first. For example, replace one Aid with a Hold Person.
\item Leave the list and choose Rest Party. Selecting spells does not recover them immediately. Wait for messages confirming preparation; an interrupted rest may leave spells unfinished.
\item Return to the adventure, use the book or symbol from a hand, choose the level and a ready spell. If a recipient is requested, select that portrait.
\item Rest again to recover the spent spell. To learn new mage magic, first scribe its scroll and then select it for preparation.
\end{enumerate}

\textbf{Preparation keyboard controls.} Up and down highlight a spell; left and right change spell level. Return or Space adds one copy, the minus key removes one, and Esc exits keeping the selection. The number on the right is the selected quantity, not a field for typing digits.

\textbf{Touch, projectile and protection.} A projectile needs a clear direction towards its target; a touch spell may require the front rank or a free hand. With Shocking Grasp, Vampiric Touch and Flame Blade you must then use the hand containing the magic. A protection does not cause damage by itself. Section 13 explains each case.

\textbf{Combat rhythm.} Attack, wait for the hand to become available and use it again. Retreating or stepping sideways can buy time, provided you have checked the floor. Do not open the inventory expecting enemies to stop. Recover projectiles left on the floor once the area is safe.
